% ═══════════════════════════════════════════════════════════════
% ML-07d: Compras - Wiederholung MUSTERLÖSUNG
% Spanisch Klasse 7 - Sequenz 7: Compras
% ═══════════════════════════════════════════════════════════════

\documentclass[11pt, a4paper]{article}

\newif\ifswmodus
\swmodusfalse

\usepackage[utf8]{inputenc}
\usepackage[T1]{fontenc}
\usepackage[ngerman]{babel}
\usepackage{helvet}
\renewcommand{\familydefault}{\sfdefault}

\usepackage[margin=2cm]{geometry}
\usepackage{enumitem}
\usepackage{tabularx}
\usepackage{booktabs}
\usepackage{multirow}

\usepackage{xcolor}
\usepackage{framed}
\usepackage{tcolorbox}
\tcbuselibrary{skins}

\usepackage{fancyhdr}
\usepackage{lastpage}
\usepackage{needspace}

% FARBEN
\definecolor{spanisch-color}{HTML}{c41e3a}
\definecolor{grau-color}{HTML}{888888}
\definecolor{hellgrau-color}{HTML}{f5f5f5}
\definecolor{loesung-color}{HTML}{c62828}

\definecolor{sw-dark}{HTML}{000000}
\definecolor{sw-gray}{HTML}{666666}
\definecolor{sw-white}{HTML}{ffffff}

\ifswmodus
    \colorlet{spanisch}{sw-dark}
    \colorlet{grau}{sw-gray}
    \colorlet{hellgrau}{sw-white}
    \colorlet{loesung}{sw-dark}
\else
    \colorlet{spanisch}{spanisch-color}
    \colorlet{grau}{grau-color}
    \colorlet{hellgrau}{hellgrau-color}
    \colorlet{loesung}{loesung-color}
\fi

\newcommand{\fachfarbe}{spanisch}

\newcommand{\thema}{Compras - Wiederholung}
\newcommand{\sequenz}{07}
\newcommand{\block}{d}
\newcommand{\fach}{Spanisch}
\newcommand{\klasse}{7}
\newcommand{\maxpunkte}{20}
\newcommand{\datum}{\today}

\pagestyle{fancy}
\fancyhf{}
\renewcommand{\headrulewidth}{0.4pt}
\renewcommand{\footrulewidth}{0.4pt}

\lhead{\fach{} -- Klasse \klasse}
\chead{\textcolor{loesung}{\textbf{ML-\sequenz\block{} MUSTERLÖSUNG}}}
\rhead{\datum}

\lfoot{}
\cfoot{}
\rfoot{Seite \thepage\ von \pageref{LastPage}}

\newcommand{\punktebox}[1]{\hfill\fbox{\small \_/#1}}
\newcommand{\luecke}[1]{\rule{#1}{0.4pt}}
\newcommand{\lsg}[1]{\textcolor{loesung}{\textbf{#1}}}

\newtcolorbox{merkkasten}[1][]{
    colback=hellgrau,
    colframe=\fachfarbe,
    colbacktitle=white,
    coltitle=\fachfarbe,
    boxrule=1pt,
    arc=0pt,
    left=8pt, right=8pt, top=8pt, bottom=8pt,
    fonttitle=\bfseries,
    attach boxed title to top left={yshift=-2mm, xshift=5mm},
    boxed title style={boxrule=0pt, colback=white},
    title=#1
}

\newenvironment{aufgabe}[1]{%
    \needspace{5\baselineskip}%
    \nopagebreak[4]%
    \vspace{10pt}
    \noindent\textbf{\textcolor{\fachfarbe}{Aufgabe #1}}
    \nopagebreak[4]%
    \vspace{5pt}
    \nopagebreak[4]%
    \begin{list}{}{\leftmargin=0pt}
    \item[]
}{%
    \end{list}
}

\newcommand{\es}[1]{\textcolor{\fachfarbe}{\textbf{#1}}}

\begin{document}

\begin{center}
    {\Large\bfseries\textcolor{\fachfarbe}{Arbeitsblatt: \thema}}\\[0.3em]
    {\large\textcolor{loesung}{\textbf{--- MUSTERLÖSUNG ---}}}
\end{center}

\vspace{5pt}

\noindent
\begin{tabularx}{\textwidth}{|X|X|X|}
\hline
\textbf{Name:} \lsg{Musterschüler/in} &
\textbf{Klasse:} \klasse &
\textbf{Datum:} \datum \\
\hline
\end{tabularx}

\vspace{15pt}

% MERKKASTEN: Lebensmittel
\begin{merkkasten}[Merkkasten: Lebensmittel (alimentos)]

\begin{tabular}{ll|ll}
\textbf{Obst (frutas):} & & \textbf{Gemüse (verduras):} & \\[0.3em]
la manzana & = \lsg{der Apfel} & el tomate & = \lsg{die Tomate} \\[0.3em]
el plátano & = \lsg{die Banane} & la lechuga & = \lsg{der Salat} \\[0.3em]
la naranja & = \lsg{die Orange} & la cebolla & = \lsg{die Zwiebel} \\[0.8em]
\textbf{Fleisch/Fisch:} & & \textbf{Milchprodukte:} & \\[0.3em]
la carne & = \lsg{das Fleisch} & la leche & = \lsg{die Milch} \\[0.3em]
el pollo & = \lsg{das Hähnchen} & el queso & = \lsg{der Käse} \\[0.3em]
el pescado & = \lsg{der Fisch} & el yogur & = \lsg{der Joghurt} \\
\end{tabular}

\end{merkkasten}

\vspace{10pt}

% AUFGABE 1
\begin{aufgabe}{1 -- Ordne zu: Lebensmittel}
Verbinde das spanische Wort mit der deutschen Bedeutung. \punktebox{4}
\end{aufgabe}

\begin{center}
\begin{tabular}{rcl}
\es{el pan} & \lsg{---} & die Milch \lsg{(= el pan $\rightarrow$ das Brot)} \\[0.8em]
\es{el agua} & \lsg{---} & das Brot \lsg{(= el agua $\rightarrow$ das Wasser)} \\[0.8em]
\es{el jamón} & \lsg{---} & das Wasser \lsg{(= el jamón $\rightarrow$ der Schinken)} \\[0.8em]
\es{la leche} & \lsg{---} & der Schinken \lsg{(= la leche $\rightarrow$ die Milch)} \\
\end{tabular}
\end{center}

\textbf{Lösung:} el pan = Brot, el agua = Wasser, el jamón = Schinken, la leche = Milch

\vspace{10pt}

% AUFGABE 2
\begin{aufgabe}{2 -- Schreibe die Zahlen als Wort}
\punktebox{4}
\end{aufgabe}

\noindent
a) 25 = \lsg{veinticinco} \hfill
b) 42 = \lsg{cuarenta y dos} \\[0.8em]
c) 78 = \lsg{setenta y ocho} \hfill
d) 100 = \lsg{cien} \\

\vspace{15pt}

% MERKKASTEN: querer/poder
\begin{merkkasten}[Merkkasten: querer (wollen) und poder (können)]

\begin{tabular}{|l|l|l|}
\hline
\textbf{Person} & \textbf{querer} & \textbf{poder} \\
\hline
yo & \lsg{quiero} & \lsg{puedo} \\
\hline
tú & \lsg{quieres} & \lsg{puedes} \\
\hline
él/ella & \lsg{quiere} & \lsg{puede} \\
\hline
nosotros & \lsg{queremos} & \lsg{podemos} \\
\hline
\end{tabular}

\vspace{0.5em}
\textbf{Achtung:} Stammvokalwechsel! e $\rightarrow$ ie (querer), o $\rightarrow$ ue (poder)
\end{merkkasten}

\vspace{10pt}

% AUFGABE 3
\begin{aufgabe}{3 -- Ergänze die richtige Form}
Wähle zwischen \es{querer} und \es{poder}. \punktebox{6}
\end{aufgabe}

\noindent
a) Yo \lsg{quiero} comprar manzanas. (Ich will Äpfel kaufen.) \\[0.8em]
b) ¿Tú \lsg{puedes} ayudarme? (Kannst du mir helfen?) \\[0.8em]
c) María \lsg{quiere} un kilo de tomates. (María will ein Kilo Tomaten.) \\[0.8em]
d) Nosotros no \lsg{podemos} pagar con tarjeta. (Wir können nicht mit Karte zahlen.) \\[0.8em]
e) ¿Vosotros \lsg{queréis} ir al mercado? (Wollt ihr auf den Markt gehen?) \\[0.8em]
f) Ellos \lsg{pueden} hablar español. (Sie können Spanisch sprechen.) \\

\vspace{15pt}

% AUFGABE 4
\begin{aufgabe}{4 -- Vervollständige den Einkaufsdialog}
Ergänze die fehlenden Teile auf Spanisch. \punktebox{6}
\end{aufgabe}

\noindent
\textbf{Verkäufer:} ¡Buenos días! ¿Qué desea?

\vspace{0.8em}

\noindent
\textbf{Kunde:} \lsg{Quiero un kilo de naranjas.} (Ich möchte ein Kilo Orangen.) \hfill (2P)

\vspace{0.8em}

\noindent
\textbf{Verkäufer:} Aquí tiene. ¿Algo más?

\vspace{0.8em}

\noindent
\textbf{Kunde:} Sí, \lsg{¿puedo tener queso?} (Kann ich auch Käse haben?) \hfill (1P)

\vspace{0.8em}

\noindent
\textbf{Verkäufer:} Claro. ¿Cuánto quiere?

\vspace{0.8em}

\noindent
\textbf{Kunde:} 200 gramos, por favor. \lsg{¿Cuánto es?} (Wie viel kostet das?) \hfill (1P)

\vspace{0.8em}

\noindent
\textbf{Verkäufer:} \lsg{Son cinco euros cincuenta.} (Das macht 5 Euro 50.) \hfill (2P)

% PUNKTEÜBERSICHT

\vspace{20pt}
\hrule
\vspace{10pt}

\begin{flushright}
\textbf{Punkte:} \lsg{20} / \maxpunkte
\end{flushright}

\end{document}
